%==============================================================================
% LSTM-uHC: Manifold-Constrained Hyper-Connections for Stable Depth Scaling
% IEEE conference format (IEEEtran). Compile: pdflatex main && bibtex main && pdflatex x2
%
% All tables in paper/tables/ are auto-generated from measured test/stability data
% by scripts/make_latex_tables.py. Figures are in ../outputs/final_ready_plots/.
%==============================================================================
\documentclass[conference]{IEEEtran}
\IEEEoverridecommandlockouts
\usepackage{cite}
\usepackage{amsmath,amssymb,amsfonts}
\usepackage{graphicx}
\usepackage{booktabs}
\usepackage{xcolor}
\usepackage{url}
\usepackage{lmodern}
\renewcommand{\ttdefault}{lmtt} % avoid missing Courier (pcr) TFM on minimal TeX
\graphicspath{{../outputs/final_ready_plots/}}

\begin{document}

\title{Manifold-Constrained Hyper-Connections for\\ Stable Depth Scaling in LSTM Time-Series Forecasting}

\author{
\IEEEauthorblockN{Chandan Kumar Shah, Abhaya Shrestha, Asmit Khanal, Kushal Regmi, Nabina Thapa}
\IEEEauthorblockA{Department of Electronics and Computer Engineering\\
Pulchowk Campus, Institute of Engineering, Tribhuvan University, Lalitpur, Nepal\\
Email: 080bct023.chandan@pcampus.edu.np, 080bct006.abhaya@pcampus.edu.np, 080bct017.asmit@pcampus.edu.np,\\
080bct042.kushal@pcampus.edu.np, 080bct047.nabina@pcampus.edu.np}}

\maketitle

\begin{abstract}
Deep stacked LSTMs hit a depth ceiling: gradients vanish or explode and deep
variants collapse to near-constant predictions. Hyper-Connections (HC) widen the
residual stream with learnable mixing matrices but are unstable; manifold-constrained
HC (mHC) restores stability for large language models via a doubly-stochastic
(Birkhoff-polytope) projection. We adapt mHC to recurrent forecasting as a
\emph{macro} inter-layer highway that keeps the Sinkhorn projection outside the
temporal loop, preserving cuDNN-optimized recurrence. On NOAA~JFK weather and
ETTh1, across depths $L\in\{4,8,16\}$ with five seeds, mHC keeps held-out test
error stable as depth increases, while Standard LSTM/TCN collapse at depth and
unconstrained HC diverges. Additional $L{=}32$ runs on weather and ETTm1 support
the same depth-stability trend but have smaller seed counts. mHC keeps the
composite residual spectral norm within numerical tolerance of $1.0$,
whereas HC's exceeds $10^{5}$. We analyze the macro design's complexity advantage
over a per-timestep (micro) placement: the Sinkhorn projection runs $O(L)$ times
rather than $O(T{\cdot}L)$. In short, the manifold constraint removes the LSTM
depth ceiling.
\end{abstract}

\begin{IEEEkeywords}
time-series forecasting, LSTM, hyper-connections, doubly stochastic matrices,
depth scaling, gradient stability.
\end{IEEEkeywords}

%==============================================================================
\section{Introduction}
\label{sec:intro}
Recurrent networks, and LSTMs~\cite{hochreiter1997lstm} in particular, remain
competitive for multivariate time-series forecasting, yet \emph{deep} stacked
LSTMs are difficult to train: as layers accumulate,
gradients vanish or explode~\cite{bengio1994difficult} and the network often collapses to a near-constant
predictor. In our experiments a 16-layer stacked LSTM on weather data degrades
from a test MSE of $92.6$ at $L{=}4$ to $204.4$ at $L{=}16$, and a temporal
convolutional network (TCN) collapses similarly. This ``depth ceiling'' limits
the representational capacity that practitioners can actually use.

Residual connections~\cite{he2016resnet,he2016identity} mitigate this through an
identity mapping that preserves signal across layers. Hyper-Connections
(HC)~\cite{zhu2024hyperconnections} generalize residuals by widening the stream
into $n$ parallel channels with learnable mixing matrices, improving expressivity
but breaking the identity mapping and, as we confirm, diverging at depth.
Manifold-constrained HC (mHC)~\cite{xie2026mhc} restores stability for large
language models by projecting the residual-mixing matrix onto the Birkhoff
polytope of doubly-stochastic matrices via the Sinkhorn--Knopp
algorithm~\cite{sinkhorn1967}, which guarantees a non-expansive, compositionally
closed mapping.

We adapt mHC to recurrent forecasting. Rather than applying the projection inside
the recurrent cell (the per-token placement used for transformers), we introduce a
\emph{macro} inter-layer highway: the doubly-stochastic mixing operates between
LSTM layers, outside the temporal loop, so the cuDNN-optimized sequence kernel is
preserved and the Sinkhorn projection runs $O(L)$ times rather than $O(T{\cdot}L)$.

\textbf{Contributions.}
\begin{itemize}
  \item A macro-level mHC highway for stacked LSTMs that decouples cross-stream
        mixing (doubly-stochastic) from temporal recurrence (cuDNN LSTM).
  \item Evidence on two five-seed datasets (NOAA~JFK weather and ETTh1) across
        $L\in\{4,8,16\}$, plus measured $L{=}32$ extensions on weather and ETTm1,
        that mHC preserves held-out test error and bounded gradients/spectral norm
        at depth while standard LSTM/TCN collapse and unconstrained HC diverges.
  \item A controlled complexity analysis showing the macro placement reduces the
        number of Sinkhorn projections from $O(T{\cdot}L)$ (per-timestep/micro) to
        $O(L)$ (per-layer), making the manifold constraint practical for long
        sequences.
\end{itemize}

%==============================================================================
\section{Related Work}
\label{sec:related}
\textbf{Residual and identity mappings.} Deep residual
networks~\cite{he2016resnet} and the identity-mapping
analysis~\cite{he2016identity} established that a clean additive skip path is
critical for trainable depth. Our work asks how to retain that property while
allowing richer cross-stream mixing.

\textbf{Hyper-Connections and mHC.} Hyper-Connections~\cite{zhu2024hyperconnections}
widen the residual stream and learn how to route information across channels, but
the unconstrained mixing matrices destroy the identity mapping and destabilize
deep training. mHC~\cite{xie2026mhc} constrains the mixing matrix to the Birkhoff
polytope via Sinkhorn--Knopp~\cite{sinkhorn1967}, demonstrated at LLM scale on the
DeepSeek-V3 architecture~\cite{liu2024deepseekv3}. The mHC paper itself frames the
choice of where to apply the highway as a \emph{macro} (inter-block) vs.\ \emph{micro}
(intra-block) design decision~\cite{xie2026mhc}. We transfer this idea to
recurrent forecasting and study \emph{where} the projection should live for an
LSTM: as a macro inter-layer highway, leaving the cuDNN-optimized sequence kernel
intact. Concurrently, mHC has also been adapted to state-space language
models~\cite{mutlu2026mhcssm}; that work extends mHC to SSMs for language
modeling, whereas we focus on LSTM-based time-series forecasting and provide a
controlled depth-scaling study.

\textbf{Time-series forecasting baselines.} Linear models such as
DLinear/NLinear~\cite{zeng2023dlinear} are strong, simple baselines; recurrent
(GRU~\cite{cho2014gru}) and convolutional (TCN~\cite{bai2018tcn}) models are
standard neural baselines; and transformer forecasters including
Informer~\cite{zhou2021informer},
PatchTST~\cite{nie2023patchtst}, iTransformer~\cite{liu2024itransformer}, and
TimesNet~\cite{wu2023timesnet} represent recent state of the art. We compare
against classical, linear, recurrent, convolutional, and (PatchTST, iTransformer)
transformer baselines under one protocol; extending to further transformer SOTA is
future work (Sec.~\ref{sec:discussion}).

\textbf{State-space models.} Selective state-space models such as
Mamba~\cite{gu2024mamba} offer an orthogonal route to stable, efficient long-sequence
modeling via input-dependent recurrences rather than a wider residual stream. We do
not train an SSM baseline in this work; mHC is a modification to the \emph{inter-layer
highway} around an existing recurrent cell (here, the LSTM) and is in principle
compatible with an SSM cell in place of the LSTM, which we leave to future work.

%==============================================================================
\section{Method}
\label{sec:method}
\textbf{Multi-stream macro highway.} We expand the residual stream into $n$
parallel streams $X_l\in\mathbb{R}^{B\times T\times n\times d}$ (we use $n{=}4$,
$d{=}128$). At each inter-layer boundary we flatten and RMS-normalize~\cite{zhang2019rmsnorm}
the streams and compute three heads from the normalized input: an aggregation gate
$\mathcal{H}^{pre}=\sigma(\cdot)\in[0,1]$, a distribution gate
$\mathcal{H}^{post}=2\sigma(\cdot)\in[0,2]$, and a residual-mixing matrix
$\mathcal{H}^{res}=\mathrm{SinkhornKnopp}(\cdot)$ projected onto the Birkhoff
polytope. Each head combines a small input-dependent (dynamic) term, gated by a
learnable $\alpha$ initialized to $0.01$, with a learned static bias, following
the mHC parameterization. The layer update is
\begin{equation}
X_{l+1}=\mathcal{H}^{res}_l X_l + \mathcal{H}^{post}_l\,
        \mathrm{LSTM}_l\!\big(\textstyle\sum_n \mathcal{H}^{pre}_l X_l\big).
\end{equation}
The aggregation gate collapses the $n$ streams into a single $d$-dimensional input
for a standard LSTM layer; the distribution gate scatters the layer output back
across streams; and the doubly-stochastic $\mathcal{H}^{res}$ mixes the streams as
a convex (soft-permutation) combination.

\textbf{Why macro placement (complexity).} Because $\mathcal{H}^{res}$ is computed
once per layer over the batched $B\times T$ tensor and the LSTM consumes the whole
sequence in one cuDNN pass, the Sinkhorn--Knopp projection executes $O(L)$ times per
forward pass. A \emph{micro} placement that applies the projection inside the
recurrent cell at every timestep would instead execute it $O(T{\cdot}L)$ times. With
$T{=}724$ this is a factor of $T$ difference ($L$ vs.\ $\sim$$724L$ projections), and
it also forgoes the cuDNN sequence kernel by requiring a step-by-step cell loop. We
therefore adopt the macro placement; Sec.~\ref{sec:micro} formalizes this argument.
We do not train a micro variant: the comparison here is analytical (operation
count), not empirical.

\subsection{Theory}
\label{sec:theory}
The Birkhoff polytope $\mathcal{M}^{res}=\{H\!\ge\!0 : H\mathbf{1}{=}\mathbf{1},
\mathbf{1}^{\!\top}H{=}\mathbf{1}^{\!\top}\}$ confers three properties that make
depth stable. \emph{(i) Non-expansiveness:} every doubly-stochastic $H$ has spectral
norm $\|H\|_2\le 1$, so the residual map cannot amplify signal energy. \emph{(ii)
Compositional closure:} the product of doubly-stochastic matrices is
doubly-stochastic, so the composite map across any number of layers
$\prod_l \mathcal{H}^{res}_l$ remains non-expansive --- stability holds at arbitrary
depth, not just per layer. \emph{(iii) Convex hull of permutations:} the polytope is
the convex hull of permutation matrices, so $\mathcal{H}^{res}$ acts as a soft
permutation that fuses streams without destructive amplification. When $n{=}1$ the
constraint reduces to the scalar $1$, recovering the standard identity-mapping
residual. Section~\ref{sec:results} confirms these properties empirically: the
measured composite spectral norm stays near $1.0$ for mHC at every depth.

%==============================================================================
\section{Experimental Setup}
\label{sec:setup}
\textbf{Datasets.} We use three real multivariate series with chronological,
gap-separated 70/15/15 train/val/test splits: the NOAA JFK hourly weather
benchmark (75{,}119 rows, 4 variables---temperature, humidity, pressure, wind
speed), ETTh1 (17{,}420 rows), and the larger ETTm1 (69{,}680 rows, 15-min
cadence). This is the NOAA JFK weather series, not the 21-variable ``Weather''
benchmark used in many long-term-forecasting papers, so all baseline numbers
reported here are from our own retraining under the protocol below. The $z$-score scaler is fit on the training split only; a guard
gap equal to the maximum horizon separates splits so no window crosses a boundary.

\textbf{Protocol.} Every model uses an identical regimen:
AdamW~\cite{loshchilov2019adamw}
($\beta{=}(0.9,0.95)$, $\epsilon{=}10^{-8}$, weight decay $0.1$), learning rate
$4\times10^{-4}$ with 5-epoch linear warmup then cosine decay, dropout $0.1$,
gradient clipping at $1.0$, batch size $32$, up to 50 epochs with early stopping
(patience 10), and five seeds (42--46) for the main weather/ETTh1 experiments.
Forecast horizons are 6/12/24/48/72 hours.
The only difference between models is the inter-layer highway. We select the
checkpoint with the best validation MSE and evaluate it \emph{once} on the held-out
test split; all reported metrics are de-normalized (original units).

\textbf{Models.} Standard LSTM (parameter-matched), unconstrained HC-LSTM, and
mHC-LSTM (ours), plus classical/linear baselines (Ridge, DLinear, NLinear) and
neural baselines (GRU, TCN), at depths $L\in\{4,8,16\}$; depth $L{=}32$ adds the
large-dataset comparison on weather and ETTm1.

%==============================================================================
\section{Results}
\label{sec:results}

\subsection{Training stability}
mHC keeps the composite residual spectral norm within numerical tolerance of $1.0$,
while unconstrained HC explodes (Table~\ref{tab:stability}, Fig.~\ref{fig:spectral})
and diverges (NaN) at $L{=}16$ on both datasets.
\begin{figure}[t]\centering
  \includegraphics[width=\linewidth]{spectral_norm_vs_depth.pdf}
  \caption{Max composite spectral norm vs.\ depth. mHC stays near $1.0$; HC explodes
  and then diverges at $L{=}16$.}\label{fig:spectral}
\end{figure}
\begin{figure}[t]\centering
  \includegraphics[width=\linewidth]{gradient_norm_vs_depth.pdf}
  \caption{Max pre-clip gradient norm vs.\ depth (log scale). mHC stays bounded
  ($<$50) at all depths; HC reaches $10^{5}$--$10^{11}$ and diverges at $L{=}16$.}\label{fig:gradnorm}
\end{figure}
\begin{table*}[t]
\centering
\caption{Training stability (max over all steps/seeds). mHC holds the composite spectral norm near $1.0$ at every depth; unconstrained HC explodes (already at $L{=}8$ on ETTh1) and diverges (NaN) at $L{=}16$.}
\label{tab:stability}
\footnotesize
\begin{tabular}{llrrr}
\toprule
Dataset & Model & $L$ & Max grad norm & Max spec. norm \\
\midrule
Weather & mHC-LSTM (ours) & 4 & 1.56e+01 & 1.0000 \\
Weather & HC-LSTM & 4 & 6.56e+05 & 628.9461 \\
Weather & Standard LSTM & 4 & 1.59e+02 & -- \\
Weather & mHC-LSTM (ours) & 8 & 9.92e+00 & 1.0000 \\
Weather & HC-LSTM & 8 & 2.47e+11 & 395483.5938 \\
Weather & Standard LSTM & 8 & 2.09e+04 & -- \\
Weather & mHC-LSTM (ours) & 16 & 6.86e+00 & 1.0001 \\
Weather & HC-LSTM & 16 & \multicolumn{2}{c}{diverged} \\
Weather & Standard LSTM & 16 & 2.38e-01 & -- \\
\midrule
ETTh1 & mHC-LSTM (ours) & 4 & 4.08e+01 & 1.0000 \\
ETTh1 & HC-LSTM & 4 & 7.22e+05 & 629.1806 \\
ETTh1 & Standard LSTM & 4 & 5.62e+02 & -- \\
ETTh1 & mHC-LSTM (ours) & 8 & 5.59e+00 & 1.0000 \\
ETTh1 & HC-LSTM & 8 & 2.57e+11 & 397134.6875 \\
ETTh1 & Standard LSTM & 8 & 5.14e+02 & -- \\
ETTh1 & mHC-LSTM (ours) & 16 & 8.50e+00 & 1.0001 \\
ETTh1 & HC-LSTM & 16 & \multicolumn{2}{c}{diverged} \\
ETTh1 & Standard LSTM & 16 & 4.14e+02 & -- \\
\bottomrule
\end{tabular}
\end{table*}


\subsection{Depth scaling (the main result)}
On held-out test data, mHC preserves stable error across the evaluated depth range;
Standard LSTM and TCN collapse at $L\geq16$
(Table~\ref{tab:depth}, Fig.~\ref{fig:depth}). The stability comes at a modest
absolute-cost increase: mHC's test MSE at $L{=}16$ on weather ($95.4$) is higher
than its own $L{=}4$ value ($91.7$), so deeper mHC does not improve raw accuracy on
these tasks; rather, it avoids the catastrophic collapse seen in the vanilla LSTM.

\textbf{Relation to stronger recurrent baselines.} We note frankly that a plain
GRU---a standard recurrent baseline with no highway machinery and fewer parameters
(351K vs.\ 592K)---achieves lower weather test MSE than mHC-LSTM at every depth
we evaluate (87.7/89.0/94.5 vs.\ 91.7/94.3/95.4 at $L{=}4/8/16$) and is
statistically tied with mHC on ETTh1 (within 0.02--0.5 MSE). Ridge and the linear
DLinear/NLinear baselines also outperform mHC on ETTh1. These results reframe the
contribution: mHC-LSTM is not a generic state-of-the-art forecaster on these
datasets at shallow depth. Its value is specific to the problem of training a
deep stacked LSTM without collapse. Practically, this means mHC-LSTM is most
interesting when an LSTM-based architecture is already preferred for deployment,
interpretability, or hardware reasons, and one wishes to add depth without the
instability that otherwise afflicts deep LSTMs.
\begin{figure}[t]\centering
  \includegraphics[width=\linewidth]{depth_scaling_test_mse_weather.pdf}
  \caption{Weather test MSE vs.\ depth. mHC stays in the 91.7--95.9 range across
  all depths; Standard LSTM and TCN collapse at $L\geq16$.}\label{fig:depth}
\end{figure}
\begin{table*}[t]
\centering
\caption{Test MSE vs.\ depth. mHC remains stable across the evaluated depth range; Standard LSTM and TCN collapse at depth.}
\label{tab:depth}
\footnotesize
\begin{tabular}{llrrrr}
\toprule
Dataset & Model & $L{=}4$ & $L{=}8$ & $L{=}16$ & $L{=}32$ \\
\midrule
Weather & mHC-LSTM (ours) & 91.7 & 94.3 & 95.4 & 95.9 \\
Weather & Standard LSTM & 92.6 & 93.1 & 204.4 & 204.6 \\
Weather & GRU & 87.7 & 89.0 & 94.5 & -- \\
Weather & TCN & 90.7 & 94.3 & 257.6 & -- \\
\midrule
ETTh1 & mHC-LSTM (ours) & 9.9 & 10.3 & 10.7 & -- \\
ETTh1 & Standard LSTM & 9.9 & 11.8 & 12.9 & -- \\
ETTh1 & GRU & 10.1 & 10.7 & 10.7 & -- \\
ETTh1 & TCN & 10.2 & 10.6 & 34.8 & -- \\
\bottomrule
\end{tabular}
\end{table*}

\begin{table*}[t]
\centering
\caption{Held-out \textbf{test}-set forecasting error (original units, mean$\pm$std over seeds). Lower is better. HC diverges (NaN) at $L{=}16$; vanilla collapses at $L{\geq}16$.}
\label{tab:main}
\footnotesize
\begin{tabular}{llrrrc}
\toprule
Dataset & Model & $L$ & Test MSE & Test MAE & $n$ \\
\midrule
Weather & mHC-LSTM (ours) & 4 & 91.69\,$\pm$\,1.41 & 5.729\,$\pm$\,0.043 & 5 \\
Weather & Standard LSTM & 4 & 92.63\,$\pm$\,1.01 & 5.823\,$\pm$\,0.037 & 5 \\
Weather & HC-LSTM & 4 & 101.81\,$\pm$\,4.86 & 6.012\,$\pm$\,0.119 & 5 \\
Weather & GRU & 4 & 87.71\,$\pm$\,0.83 & 5.609\,$\pm$\,0.035 & 5 \\
Weather & TCN & 4 & 90.68\,$\pm$\,1.13 & 5.722\,$\pm$\,0.042 & 5 \\
Weather & Ridge & 4 & 88.87\,$\pm$\,1.21 & 5.644\,$\pm$\,0.053 & 5 \\
Weather & DLinear & 4 & 98.22\,$\pm$\,1.03 & 6.056\,$\pm$\,0.054 & 4 \\
Weather & NLinear & 4 & 96.69\,$\pm$\,0.74 & 5.939\,$\pm$\,0.013 & 5 \\
Weather & mHC-LSTM (ours) & 8 & 94.29\,$\pm$\,0.85 & 5.820\,$\pm$\,0.025 & 5 \\
Weather & Standard LSTM & 8 & 93.12\,$\pm$\,2.30 & 5.883\,$\pm$\,0.087 & 5 \\
Weather & HC-LSTM & 8 & 94.00\,$\pm$\,2.71 & 5.795\,$\pm$\,0.094 & 5 \\
Weather & GRU & 8 & 88.99\,$\pm$\,1.19 & 5.672\,$\pm$\,0.066 & 5 \\
Weather & TCN & 8 & 94.29\,$\pm$\,0.98 & 5.909\,$\pm$\,0.044 & 5 \\
Weather & mHC-LSTM (ours) & 16 & 95.37\,$\pm$\,1.39 & 5.863\,$\pm$\,0.040 & 5 \\
Weather & Standard LSTM & 16 & 204.38\,$\pm$\,0.29 & 9.671\,$\pm$\,0.005 & 5 \\
Weather & GRU & 16 & 94.53\,$\pm$\,3.04 & 5.895\,$\pm$\,0.108 & 5 \\
Weather & TCN & 16 & 257.63\,$\pm$\,9.31 & 10.583\,$\pm$\,0.156 & 4 \\
Weather & mHC-LSTM (ours) & 32 & 95.89\,$\pm$\,2.44 & 5.878\,$\pm$\,0.061 & 5 \\
Weather & Standard LSTM & 32 & 204.55\,$\pm$\,0.39 & 9.672\,$\pm$\,0.009 & 3 \\
\midrule
ETTh1 & mHC-LSTM (ours) & 4 & 9.92\,$\pm$\,0.26 & 1.952\,$\pm$\,0.046 & 5 \\
ETTh1 & Standard LSTM & 4 & 9.92\,$\pm$\,0.71 & 2.028\,$\pm$\,0.089 & 5 \\
ETTh1 & HC-LSTM & 4 & 44.52\,$\pm$\,18.73 & 4.323\,$\pm$\,0.983 & 5 \\
ETTh1 & GRU & 4 & 10.08\,$\pm$\,0.22 & 2.014\,$\pm$\,0.046 & 5 \\
ETTh1 & TCN & 4 & 10.16\,$\pm$\,0.44 & 2.097\,$\pm$\,0.109 & 5 \\
ETTh1 & Ridge & 4 & 7.77\,$\pm$\,0.09 & 1.691\,$\pm$\,0.020 & 5 \\
ETTh1 & DLinear & 4 & 8.35\,$\pm$\,0.06 & 1.718\,$\pm$\,0.011 & 5 \\
ETTh1 & NLinear & 4 & 8.26\,$\pm$\,0.08 & 1.712\,$\pm$\,0.017 & 5 \\
ETTh1 & mHC-LSTM (ours) & 8 & 10.26\,$\pm$\,0.56 & 1.995\,$\pm$\,0.129 & 5 \\
ETTh1 & Standard LSTM & 8 & 11.79\,$\pm$\,3.06 & 2.254\,$\pm$\,0.351 & 5 \\
ETTh1 & HC-LSTM & 8 & 8605623.25\,$\pm$\,3713064.63 & 1900.638\,$\pm$\,334.541 & 5 \\
ETTh1 & GRU & 8 & 10.74\,$\pm$\,0.36 & 2.048\,$\pm$\,0.049 & 5 \\
ETTh1 & TCN & 8 & 10.57\,$\pm$\,1.17 & 2.153\,$\pm$\,0.110 & 5 \\
ETTh1 & mHC-LSTM (ours) & 16 & 10.65\,$\pm$\,0.66 & 2.058\,$\pm$\,0.108 & 5 \\
ETTh1 & Standard LSTM & 16 & 12.88\,$\pm$\,2.60 & 2.378\,$\pm$\,0.292 & 5 \\
ETTh1 & GRU & 16 & 10.67\,$\pm$\,0.48 & 2.116\,$\pm$\,0.055 & 5 \\
ETTh1 & TCN & 16 & 34.84\,$\pm$\,4.88 & 4.168\,$\pm$\,0.369 & 5 \\
\midrule
ETTm1 & mHC-LSTM (ours) & 32 & 9.13\,$\pm$\,0.40 & 1.778\,$\pm$\,0.037 & 5 \\
ETTm1 & Standard LSTM & 32 & 37.57\,$\pm$\,0.05 & 4.322\,$\pm$\,0.005 & 4 \\
\bottomrule
\end{tabular}
\end{table*}

\begin{table*}[t]
\centering
\caption{Paired significance on \textbf{test} MSE for matched five-seed comparisons: Cohen's $d$ and two-sided $t$-test $p$. Negative diff favors mHC.}
\label{tab:sig}
\footnotesize
\begin{tabular}{llrrrc}
\toprule
Dataset & Comparison & $L$ & $\Delta$MSE & Cohen's $d$ & $p$ \\
\midrule
ETTh1 & mHC vs HC & 4 & -34.60 & -1.84 & 0.015$^{*}$ \\
ETTh1 & mHC vs Vanilla & 4 & -0.01 & -0.01 & 0.989 \\
ETTh1 & mHC vs HC & 8 & -8605612.99 & -2.32 & 0.007$^{*}$ \\
ETTh1 & mHC vs Vanilla & 8 & -1.53 & -0.51 & 0.314 \\
ETTh1 & mHC vs Vanilla & 16 & -2.23 & -0.70 & 0.192 \\
Weather & mHC vs HC & 4 & -10.12 & -1.74 & 0.018$^{*}$ \\
Weather & mHC vs Vanilla & 4 & -0.94 & -0.51 & 0.318 \\
Weather & mHC vs HC & 8 & +0.29 & 0.08 & 0.860 \\
Weather & mHC vs Vanilla & 8 & +1.17 & 0.41 & 0.406 \\
Weather & mHC vs Vanilla & 16 & -109.01 & -76.72 & 0.000$^{*}$ \\
\bottomrule
\end{tabular}
\\[2pt]
\footnotesize $^{*}$ significant at $p<0.05$.
\end{table*}


\subsection{Larger dataset (ETTm1)}
On ETTm1 (69{,}680 rows) at $L{=}32$, mHC achieves $9.13{\pm}0.40$ test MSE
(5 seeds) while Standard LSTM degrades to $37.57{\pm}0.05$ (4 seeds; one vanilla
seed failed to complete and is excluded, not imputed),
a $4.1{\times}$ gap (Fig.~\ref{fig:ettm1}).
\begin{figure}[t]\centering
  \includegraphics[width=\linewidth]{ettm1_l32_mhc_vs_vanilla.pdf}
  \caption{ETTm1 $L{=}32$: mHC scales; Standard LSTM degrades.}\label{fig:ettm1}
\end{figure}

\subsection{Baseline comparison and horizon degradation}
At shallow depth ($L{=}4$) mHC is competitive with, but not statistically distinct
from, the strongest classical, linear, recurrent, and convolutional baselines
(Fig.~\ref{fig:leaderboard}); the benefit of mHC appears with depth rather than at
$L{=}4$. Per-horizon test MAE degrades gracefully for mHC across the 6--72\,h
horizons, whereas collapsed deep baselines flatten to a near-constant error
(Fig.~\ref{fig:horizon}).
\begin{figure}[t]\centering
  \includegraphics[width=\linewidth]{test_leaderboard_L4_weather.pdf}
  \caption{Shallow ($L{=}4$) weather test MSE across all model families.}\label{fig:leaderboard}
\end{figure}
\begin{figure}[t]\centering
  \includegraphics[width=\linewidth]{horizon_mae_weather_l16.pdf}
  \caption{Per-horizon test MAE at $L{=}16$ on weather. Standard LSTM flattens
  (near-constant prediction) while mHC degrades gracefully with horizon.}\label{fig:horizon}
\end{figure}

\subsection{Comparison with modern forecasters}
\label{sec:sota}
We re-train three strong recent forecasters --- PatchTST, iTransformer, and
TimesNet --- under our \emph{identical} protocol (same splits, scaler,
horizons, optimizer, 5 seeds, checkpoint selection), so the comparison is on our
data, not copied numbers (Table~\ref{tab:sota}). On \textbf{weather}, mHC ($91.7$)
significantly outperforms both PatchTST ($95.9$, $d{=}{-}3.5$, $p{=}0.001$) and
iTransformer ($98.5$, $d{=}{-}2.3$, $p{=}0.007$) on test MSE
(Table~\ref{tab:sotasig}); TimesNet ($96.2$) is also significantly worse than
mHC on this dataset ($d{=}{-}1.36$, $p{=}0.038$), though with higher seed variance
($\sigma{=}3.3$). Our TimesNet baseline is reimplemented following the module
structure of the Time-Series-Library benchmark~\cite{wang2024tssurvey}.
On \textbf{ETTh1}, PatchTST ($8.5$) significantly beats
mHC ($9.9$), while mHC significantly beats iTransformer ($10.7$); TimesNet ($9.4$)
is non-significantly better ($p{=}0.348$).
Thus mHC is competitive with 2023--2024 SOTA --- leading on weather, mid-pack on ETTh1 ---
which, combined with its unique depth-stability (Sec.~\ref{sec:results}A--B),
supports the contribution without overclaiming shallow accuracy.
\begin{table*}[t]
\centering
\caption{Comparison against modern forecasters (held-out \textbf{test} MSE/MAE, mean$\pm$std over five seeds). Our LSTM models are at $L{=}4$; transformer baselines use their own depth. mHC is best on weather; PatchTST/NLinear lead on ETTh1.}
\label{tab:sota}
\footnotesize
\begin{tabular}{lrrrr}
\toprule
Model & ETTh1 MSE & ETTh1 MAE & Weather MSE & Weather MAE \\
\midrule
NLinear & 8.26\,$\pm$\,0.08 & 1.712\,$\pm$\,0.017 & 96.69\,$\pm$\,0.74 & 5.939\,$\pm$\,0.013 \\
DLinear & 8.35\,$\pm$\,0.06 & 1.718\,$\pm$\,0.011 & 98.22\,$\pm$\,1.03 & 6.056\,$\pm$\,0.054 \\
Ridge & 7.77\,$\pm$\,0.09 & 1.691\,$\pm$\,0.020 & 88.87\,$\pm$\,1.21 & 5.644\,$\pm$\,0.053 \\
TCN & 10.16\,$\pm$\,0.44 & 2.097\,$\pm$\,0.109 & 90.68\,$\pm$\,1.13 & 5.722\,$\pm$\,0.042 \\
GRU & 10.08\,$\pm$\,0.22 & 2.014\,$\pm$\,0.046 & 87.71\,$\pm$\,0.83 & 5.609\,$\pm$\,0.035 \\
PatchTST & 8.51\,$\pm$\,0.34 & 1.833\,$\pm$\,0.050 & 95.92\,$\pm$\,1.05 & 5.971\,$\pm$\,0.032 \\
iTransformer & 10.74\,$\pm$\,0.52 & 2.065\,$\pm$\,0.049 & 98.47\,$\pm$\,1.72 & 6.116\,$\pm$\,0.059 \\
TimesNet & 9.40\,$\pm$\,0.86 & 1.998\,$\pm$\,0.144 & 96.25\,$\pm$\,3.34 & 5.932\,$\pm$\,0.116 \\
\textbf{mHC-LSTM (ours)} & 9.92\,$\pm$\,0.26 & 1.952\,$\pm$\,0.046 & 91.69\,$\pm$\,1.41 & 5.729\,$\pm$\,0.043 \\
\bottomrule
\end{tabular}
\end{table*}

\begin{table*}[t]
\centering
\caption{Paired significance of mHC ($L{=}4$) vs.\ SOTA on \textbf{test} MSE (five matched seeds). Negative diff favors mHC.}
\label{tab:sotasig}
\footnotesize
\begin{tabular}{llrrc}
\toprule
Dataset & vs.\ & $\Delta$MSE & Cohen's $d$ & $p$ \\
\midrule
Weather & PatchTST & -4.23 & -3.52 & 0.001$^{*}$ \\
Weather & iTransformer & -6.79 & -2.30 & 0.007$^{*}$ \\
Weather & TimesNet & -4.56 & -1.36 & 0.038$^{*}$ \\
ETTh1 & PatchTST & +1.41 & 2.54 & 0.005$^{*}$ \\
ETTh1 & iTransformer & -0.82 & -1.76 & 0.017$^{*}$ \\
ETTh1 & TimesNet & +0.52 & 0.48 & 0.348 \\
\bottomrule
\end{tabular}
\\[2pt]
\footnotesize $^{*}$ significant at $p<0.05$.
\end{table*}


\subsection{Compute efficiency}
\label{sec:efficiency}
mHC adds negligible parameter/FLOP overhead over a standard LSTM: at every depth
the two are parameter- and FLOP-matched (e.g., at $L{=}4$, $592$K vs.\ $590$K parameters and
$0.85$ vs.\ $0.86$ GFLOPs), since the $n{=}4$ stream heads are tiny relative to the
LSTM matmuls (Table~\ref{tab:efficiency}). The cost of the multi-stream highway is
\emph{activation memory}: mHC's peak VRAM is higher than the standard LSTM because
it carries $n$ streams, and \emph{CPU inference latency}: mHC ($1539$ms) is $\sim$3$\times$
slower than the standard LSTM ($508$ms) due to sequential stream mixing, while
iTransformer ($3.4$ms) benefits from parallel attention in this CPU benchmark.
Parameters and FLOPs are measured on a $1{\times}724{\times}4$ input; latency is
mean over 100 CPU forward passes (batch=32). We do not claim GPU inference latency
from this CPU measurement.
\begin{table*}[t]
\centering
\caption{Compute efficiency (weather). Params and FLOPs measured on a $1{\times}724{\times}4$ input; peak VRAM and train time from training logs; inference latency on CPU (batch=32, 100 runs, mean). mHC is parameter- and FLOP-matched to the standard LSTM, at higher activation memory and CPU latency from the $n$-stream expansion.}
\label{tab:efficiency}
\footnotesize
\begin{tabular}{llrrrrr}
\toprule
Model & $L$ & Params & GFLOPs & VRAM (MB) & Train/epoch (s) & Latency (ms) \\
\midrule
mHC-LSTM (ours) & 4 & 592,512 & 0.845 & 1539 & 182.8 & 1539.3 \\
Standard LSTM & 4 & 590,420 & 0.857 & 636 & 296.3 & 507.8 \\
HC-LSTM & 4 & 592,512 & 0.845 & 1249 & 149.3 & 1349.9 \\
GRU & 4 & 351,252 & 0.510 & -- & -- & 470.7 \\
TCN & 4 & 153,108 & 0.216 & -- & -- & 131.6 \\
PatchTST & 4 & 446,725 & 0.098 & -- & -- & 106.8 \\
iTransformer & 4 & 490,885 & 0.002 & -- & -- & 3.4 \\
TimesNet & 4 & 4,738,157 & 13.572 & -- & -- & 5457.9 \\
DLinear & 4 & 7,250 & 0.000 & -- & -- & 2.0 \\
NLinear & 4 & 3,625 & 0.000 & -- & -- & 0.3 \\
\bottomrule
\end{tabular}
\end{table*}


\subsection{Macro vs.\ micro placement: a complexity argument}
\label{sec:micro}
The doubly-stochastic projection can be placed either between layers (\emph{macro},
ours) or inside the recurrent cell at every timestep (\emph{micro}). We compare the
two analytically. Let $K$ be the Sinkhorn--Knopp iteration count and $n$ the number
of streams. Per forward pass over a length-$T$ sequence with $L$ layers, the macro
placement performs $L$ projections (each a batched $n{\times}n$ Sinkhorn over the
$B{\times}T$ elements), i.e.\ $O(L\,K\,n^2)$ Sinkhorn work \emph{outside} the
temporal recurrence, so the LSTM still runs as a single fused cuDNN kernel. The micro
placement performs one projection \emph{per timestep per layer}, i.e.\ $O(T\,L\,K\,n^2)$,
and forces a step-by-step \texttt{LSTMCell} loop that forgoes the cuDNN sequence
kernel. For our setting ($T{=}724$) the micro placement issues $\sim$$724\times$ more
projections and loses hardware-optimized recurrence. We therefore adopt the macro
placement on complexity grounds; we do not train a micro variant, so no empirical
accuracy comparison is claimed.

%==============================================================================
\section{Discussion and Limitations}
\label{sec:discussion}
\textbf{What the results support.} The central claim --- that doubly-stochastic
inter-layer mixing yields stable, trainable depth for LSTMs --- is supported on
held-out test data: mHC's error remains stable with depth, its composite
spectral norm stays near $1.0$, and its gradient norm stays bounded, whereas standard
LSTM and TCN collapse at $L\!\geq\!16$ and unconstrained HC diverges. The advantage
over HC is statistically significant at $L{=}4$ on both datasets, and the advantage
over standard LSTM at $L{=}16$ on weather is large ($d\!\approx\!-77$, $p<10^{-3}$).

\textbf{What we do not claim.} At \emph{shallow} depth mHC is statistically
indistinguishable from a standard LSTM and from the strongest baselines on test
data; this is a stability/trainability result, not a claim of shallow
state-of-the-art accuracy. \textbf{Statistical caveats.} All significance tests use five matched seeds.
Large Cohen's $d$ values (e.g., $d{\approx}{-}77$ for weather $L{=}16$) reflect a
large mean gap relative to the very small seed-to-seed variance in that paired
comparison, but at $n{=}5$ such point estimates are fragile. We report uncorrected
paired $t$-test $p$-values across roughly seventeen comparisons in
Tables~\ref{tab:sig} and~\ref{tab:sotasig}; a few borderline claims
($p{\approx}0.038$) would not survive a Holm--Bonferroni multiple-comparison
correction. The robust conclusions are the qualitative stability pattern
(spectral norm near 1.0, bounded gradients) and the large weather depth-collapse
gap, not fine-grained ranking of methods within 1--2 MSE points.

\textbf{Limitations.} The ETTm1 result has 5 seeds for
mHC and 4 for vanilla (one vanilla seed failed to complete and is excluded, not
imputed). Our SOTA comparison covers PatchTST,
iTransformer, and TimesNet (5 seeds each); extending to FEDformer and Autoformer,
and to additional long-term-forecasting benchmarks, is the clearest next step. Two
baseline cells within the main five-seed sweep have only four seeds each
(weather DLinear $L{=}4$, weather TCN $L{=}16$; one seed failed to produce a
log in each case), reported as $n{=}4$ in Table~\ref{tab:main} rather than
imputed. We built and unit-tested a micro (per-timestep) mHC-LSTM and staged a
15-run training ablation (3 depths $\times$ 5 seeds); these runs were not
completed, so Sec.~\ref{sec:micro}'s comparison remains analytical only, with
no empirical micro-model accuracy numbers reported anywhere in this paper.
All models share the same learning-rate schedule across depths; we did not
perform a per-depth learning-rate search for the Standard LSTM, so part of the
depth-collapse effect could in principle be a tuning artifact rather than a
fundamental depth limit. Re-tuning depth-specific LRs is an open ablation.

\section{Conclusion}
\label{sec:conclusion}
We presented a macro-level manifold-constrained Hyper-Connection highway for stacked
LSTMs that keeps cross-stream mixing doubly-stochastic while preserving
cuDNN-optimized recurrence. Across measured real-data experiments up to $L{=}32$,
it trains stably and avoids the depth-collapse that afflicts standard LSTMs, TCNs,
and unconstrained Hyper-Connections. In short, the
manifold constraint removes the LSTM depth ceiling.

%==============================================================================
\section*{Reproducibility}
Metrics are de-normalized test-set values. The main weather/ETTh1 experiments use
five seeds; smaller-$n$ extensions are marked in the tables. Checkpoints are selected
by best validation MSE and evaluated once on the held-out test split. Tables are
generated directly from measured logs; missing/diverged runs are reported, never
imputed. Code, configs, the training log manifest, and the `make\_latex\_tables.py`
script will be released publicly upon acceptance; until then they are available
from the authors on request.

\bibliographystyle{IEEEtran}
\bibliography{refs}

\end{document}
